% GENERATED FILE. Edit canonical lessons, then run npm run generate:books.
% canonical-source-hash: fnv1a64:bd2f6b8e3da7522b
% canonical-lessons: RU-C508-bolelshchik, RU-C508-tribuna, RU-C508-svistok, RU-C508-podbadrivat, RU-C508-likovat, RU-R508-review-501-508

\chapter{In the Stands}
\label{ch:ru-in-the-stands}
\hlchaptermodality{508}

\begin{chapteropening}
\textbf{By the end of this chapter:} \emph{I can talk about fans cheering at a match.}

Complete the last lesson of chapter 508: I can talk about fans cheering at a match.
\end{chapteropening}

\section[\texorpdfstring{\ru{болельщик} (bolélshchik)}{болельщик (bolélshchik)}]{\ru{болельщик} (bolélshchik) — a fan, a supporter}

\label{lesson:RU-C508-bolelshchik}



\begin{quote}
Before the new one: say the Russian for a greenhouse, then the Russian for to prune.
\end{quote}

\subsection*{You'll want to know: \ru{болельщик}}
\textbf{\ru{болельщик}} (\emph{bolélshchik}) — \textquotedblleft{}a fan, a supporter\textquotedblright{}.

\begin{quote}
\ru{Болельщики} \ru{пели} \ru{весь} \ru{матч}.
\end{quote}

\begin{quote}
\emph{Bolélshchiki péli ves match.}
\end{quote}

\begin{quote}
The fans sang throughout the match.
\end{quote}

Masculine; feminine \ru{болельщица}. \emph{\ru{Болеть} \ru{за}} + accusative — to support a team.

\begin{cousinweb}
\emph{\ru{болельщик}} is built on \emph{\ru{болеть}} 'to be ill', and so 'to suffer for, to support', with the ending \emph{-\ru{льщик}} (EtymDB).
\end{cousinweb}

\begin{grammarlens}[title={What you've built}]
One more word for telling, explaining and describing at B1.
\end{grammarlens}

\subsection*{Your turn}
\begin{itemize}
\raggedright
  \item \emph{Say it:} \emph{bolélshchik}
  \item \emph{Say it:} \emph{bolélshchik}, once more
  \item \emph{Say it:} a sentence of your own with \emph{bolélshchik}
  \item \emph{Recall:} say the Russian for a greenhouse, then the Russian for to prune, then say \emph{bolélshchik} again
\end{itemize}

\begin{tcolorbox}[breakable,colback=teal!4,colframe=teal!35!black,title={Before you move on}]
What does \ru{болельщик} mean? (\textquotedblleft{}A fan, a supporter\textquotedblright{}.) Say it once more.
\end{tcolorbox}
\section[\texorpdfstring{\ru{трибуна} (tribúna)}{трибуна (tribúna)}]{\ru{трибуна} (tribúna) — a stand (in a stadium)}

\label{lesson:RU-C508-tribuna}



\begin{quote}
Before the new one: say the Russian for to prune, then the Russian for a fan.
\end{quote}

\subsection*{You'll want to know: \ru{трибуна}}
\textbf{\ru{трибуна}} (\emph{tribúna}) — \textquotedblleft{}a stand (in a stadium)\textquotedblright{}.

\begin{quote}
\ru{Все} \ru{трибуны} \ru{были} \ru{полны}.
\end{quote}

\begin{quote}
\emph{Vse tribúny býli polný.}
\end{quote}

\begin{quote}
All the stands were full.
\end{quote}

Feminine. Also a speaker's platform.

\begin{cousinweb}
\emph{\ru{трибуна}} goes back to Latin \emph{tribūnus} 'a tribune' (EtymDB), usually through Italian \emph{tribuna}, a raised platform.
\end{cousinweb}

\begin{grammarlens}[title={What you've built}]
One more word for telling, explaining and describing at B1.
\end{grammarlens}

\subsection*{Your turn}
\begin{itemize}
\raggedright
  \item \emph{Say it:} \emph{tribúna}
  \item \emph{Say it:} \emph{tribúna}, once more
  \item \emph{Say it:} a sentence of your own with \emph{tribúna}
  \item \emph{Recall:} say the Russian for to prune, then the Russian for a fan, then say \emph{tribúna} again
\end{itemize}

\begin{tcolorbox}[breakable,colback=teal!4,colframe=teal!35!black,title={Before you move on}]
What does \ru{трибуна} mean? (\textquotedblleft{}A stand (in a stadium)\textquotedblright{}.) Say it once more.
\end{tcolorbox}
\section[\texorpdfstring{\ru{свисток} (svistók)}{свисток (svistók)}]{\ru{свисток} (svistók) — a whistle}

\label{lesson:RU-C508-svistok}



\begin{quote}
Before the new one: say the Russian for a fan, then the Russian for a stand.
\end{quote}

\subsection*{You'll want to know: \ru{свисток}}
\textbf{\ru{свисток}} (\emph{svistók}) — \textquotedblleft{}a whistle\textquotedblright{}.

\begin{quote}
\ru{После} \ru{свистка} \ru{судьи} \ru{игра} \ru{закончилась}.
\end{quote}

\begin{quote}
\emph{Pósle svistká sudí igrá zakónchilas.}
\end{quote}

\begin{quote}
After the referee's whistle the game ended.
\end{quote}

Masculine; genitive \ru{свистка}.

\begin{cousinweb}
\emph{\ru{свисток}} is built on \emph{\ru{свистеть}} 'to whistle', an imitative root; EtymDB lists \emph{\ru{свисток}} without a source.
\end{cousinweb}

\begin{grammarlens}[title={What you've built}]
One more word for telling, explaining and describing at B1.
\end{grammarlens}

\subsection*{Your turn}
\begin{itemize}
\raggedright
  \item \emph{Say it:} \emph{svistók}
  \item \emph{Say it:} \emph{svistók}, once more
  \item \emph{Say it:} a sentence of your own with \emph{svistók}
  \item \emph{Recall:} say the Russian for a fan, then the Russian for a stand, then say \emph{svistók} again
\end{itemize}

\begin{tcolorbox}[breakable,colback=teal!4,colframe=teal!35!black,title={Before you move on}]
What does \ru{свисток} mean? (\textquotedblleft{}A whistle\textquotedblright{}.) Say it once more.
\end{tcolorbox}
\section[\texorpdfstring{\ru{подбадривать} (podbádrivat)}{подбадривать (podbádrivat)}]{\ru{подбадривать} (podbádrivat) — to cheer on, to encourage}

\label{lesson:RU-C508-podbadrivat}



\begin{quote}
Before the new one: say the Russian for a stand, then the Russian for a whistle.
\end{quote}

\subsection*{You'll want to know: \ru{подбадривать}}
\textbf{\ru{подбадривать}} (\emph{podbádrivat}) — \textquotedblleft{}to cheer on, to encourage\textquotedblright{}.

\begin{quote}
\ru{Родители} \ru{подбадривают} \ru{детей} \ru{во} \ru{время} \ru{соревнований}.
\end{quote}

\begin{quote}
\emph{Rodíteli podbádrivayut detéy vo vrémya sorevnovániy.}
\end{quote}

\begin{quote}
Parents cheer their children on during competitions.
\end{quote}

Imperfective; perfective \ru{подбодрить}. Built on \emph{\ru{бодрый}} — brisk, cheerful.

\begin{cousinweb}
\emph{\ru{подбадривать}} is built on \emph{\ru{под}-} and \emph{\ru{бодрый}} 'brisk, cheerful'; EtymDB has no entry for the verb, and \emph{\ru{бодрый}} is inherited from Proto-Slavic \emph{b\ru{ъ}dr\ru{ъ}} 'alert, cheerful' (EtymDB).
\end{cousinweb}

\begin{grammarlens}[title={What you've built}]
One more word for telling, explaining and describing at B1.
\end{grammarlens}

\subsection*{Your turn}
\begin{itemize}
\raggedright
  \item \emph{Say it:} \emph{podbádrivat}
  \item \emph{Say it:} \emph{podbádrivat}, once more
  \item \emph{Say it:} a sentence of your own with \emph{podbádrivat}
  \item \emph{Recall:} say the Russian for a stand, then the Russian for a whistle, then say \emph{podbádrivat} again
\end{itemize}

\begin{tcolorbox}[breakable,colback=teal!4,colframe=teal!35!black,title={Before you move on}]
What does \ru{подбадривать} mean? (\textquotedblleft{}To cheer on, to encourage\textquotedblright{}.) Say it once more.
\end{tcolorbox}
\section[\texorpdfstring{\ru{ликовать} (likovát)}{ликовать (likovát)}]{\ru{ликовать} (likovát) — to rejoice}

\label{lesson:RU-C508-likovat}



\begin{quote}
Before the new one: say the Russian for a whistle, then the Russian for to cheer on.
\end{quote}

\subsection*{You'll want to know: \ru{ликовать}}
\textbf{\ru{ликовать}} (\emph{likovát}) — \textquotedblleft{}to rejoice\textquotedblright{}.

\begin{quote}
\ru{После} \ru{победы} \ru{весь} \ru{стадион} \ru{ликовал}.
\end{quote}

\begin{quote}
\emph{Pósle pobédy ves stadión likovál.}
\end{quote}

\begin{quote}
After the victory the whole stadium rejoiced.
\end{quote}

Imperfective; present \ru{ликую}, \ru{ликуешь}. Rather bookish.

\begin{cousinweb}
\emph{\ru{ликовать}} is usually traced to Church Slavonic \emph{\ru{ликъ}} 'a chorus, a round dance', so rejoicing was first dancing; EtymDB lists the verb without a source.
\end{cousinweb}

\begin{grammarlens}[title={What you've built}]
One more word for telling, explaining and describing at B1.
\end{grammarlens}

\subsection*{Your turn}
\begin{itemize}
\raggedright
  \item \emph{Say it:} \emph{likovát}
  \item \emph{Say it:} \emph{likovát}, once more
  \item \emph{Say it:} a sentence of your own with \emph{likovát}
  \item \emph{Recall:} say the Russian for a whistle, then the Russian for to cheer on, then say \emph{likovát} again
\end{itemize}

\begin{tcolorbox}[breakable,colback=teal!4,colframe=teal!35!black,title={Before you move on}]
What does \ru{ликовать} mean? (\textquotedblleft{}To rejoice\textquotedblright{}.) Say it once more.
\end{tcolorbox}
\section[(dialogue)]{Words from chapters 501–508 — a review}

\label{lesson:RU-R508-review-501-508}



\begin{quote}
Say the Russian for a thread and a needle.
\end{quote}

\subsection*{Your turn}
Say these aloud:

\begin{itemize}
\raggedright
  \item a thread, a needle, a thimble, a sewing pattern, to hem
  \item clay, modelling clay, a potter, to make things by hand, a handmade craft item
  \item to collect, a collector, a badge, a crossword, to solve
  \item draughts, a pawn, a die, a pack of cards, to cheat
  \item fishing, a fishing rod, a hook, to bite, a catch
  \item a garden bed, a watering can, to weed, a weed, a harvest
  \item seedlings, to bloom, to wilt, a greenhouse, to prune
  \item a fan, a stand, a whistle, to cheer on, to rejoice
\end{itemize}

\begin{tcolorbox}[breakable,colback=teal!4,colframe=teal!35!black,title={Before you move on}]
A thread? (\textbf{\ru{нитка}}.) A needle? (\textbf{\ru{иголка}}.) A thimble? (\textbf{\ru{напёрсток}}.) A sewing pattern? (\textbf{\ru{выкройка}}.) To hem? (\textbf{\ru{подшивать}}.) Clay? (\textbf{\ru{глина}}.) Modelling clay? (\textbf{\ru{пластилин}}.) A potter? (\textbf{\ru{гончар}}.) To make things by hand? (\textbf{\ru{мастерить}}.) A handmade craft item? (\textbf{\ru{поделка}}.) To collect? (\textbf{\ru{коллекционировать}}.) A collector? (\textbf{\ru{коллекционер}}.) A badge? (\textbf{\ru{значок}}.) A crossword? (\textbf{\ru{кроссворд}}.) To solve? (\textbf{\ru{разгадывать}}.) Draughts? (\textbf{\ru{шашки}}.) A pawn? (\textbf{\ru{пешка}}.) A die? (\textbf{\ru{кубик}}.) A pack of cards? (\textbf{\ru{колода}}.) To cheat? (\textbf{\ru{жульничать}}.) Fishing? (\textbf{\ru{рыбалка}}.) A fishing rod? (\textbf{\ru{удочка}}.) A hook? (\textbf{\ru{крючок}}.) To bite? (\textbf{\ru{клевать}}.) A catch? (\textbf{\ru{улов}}.) A garden bed? (\textbf{\ru{грядка}}.) A watering can? (\textbf{\ru{лейка}}.) To weed? (\textbf{\ru{полоть}}.) A weed? (\textbf{\ru{сорняк}}.) A harvest? (\textbf{\ru{урожай}}.) Seedlings? (\textbf{\ru{рассада}}.) To bloom? (\textbf{\ru{цвести}}.) To wilt? (\textbf{\ru{вянуть}}.) A greenhouse? (\textbf{\ru{теплица}}.) To prune? (\textbf{\ru{подрезать}}.) A fan? (\textbf{\ru{болельщик}}.) A stand? (\textbf{\ru{трибуна}}.) A whistle? (\textbf{\ru{свисток}}.) To cheer on? (\textbf{\ru{подбадривать}}.) To rejoice? (\textbf{\ru{ликовать}}.)
\end{tcolorbox}
